\chapter*{Lista symboli}

\begin{itemize}[noitemsep,topsep=0pt,parsep=0pt,partopsep=0pt,labelwidth=2cm,leftmargin=2.3cm,align=left,itemindent=0pt]
\item[\(N\)] - liczba próbek transformaty
\item[\(n\)] - indeks próbki w dziedzinie czasu
\item[\(k\)] - indeks prążka widma
\item[\(x(n)\)] - \textit{n}-ta próbka sygnału dyskretnego
\item[\(X(k)\)] - \textit{k}-ty prążek widma dyskretnego
\item[\(W_N\)] - czynnik obrotu transformaty o długości \(N\)
\item[\(p,\, q\)] - pozycje dwóch elementów przetwarzanych w operacji motylkowej
\item[\(s\)] - numer etapu algorytmu FFT
\item[\(t_{1/2}\)] - mediana czasu wykonania transformaty
\item[\(Q_1,\, Q_3\)] - pierwszy i trzeci kwartyl czasu wykonania
\item[\(t_{\mathrm{CPU}}\)] - sumaryczny czas procesora
\item[\(t_{\mathrm{wall}}\)] - czas ściany, czyli rzeczywisty czas trwania obliczeń
\item[\(R\)] - stosunek czasu procesora do czasu ściany
\item[\(I\)] - liczba wykonanych instrukcji
\item[\(C\)] - liczba cykli procesora
\item[\(m\)] - współczynnik chybień danego poziomu pamięci podręcznej
\item[\(r,\, a\)] - liczba chybień i liczba odwołań do danego poziomu pamięci podręcznej
\item[\(v\)] - udział instrukcji wektorowych w kodzie wynikowym
\item[\(n_{\mathrm{SIMD}},\, n_{\mathrm{skal}}\)] - liczba instrukcji wektorowych i skalarnych
\item[\(X_k,\, X_k^{\mathrm{ref}}\)] - \textit{k}-ty prążek widma badanego wariantu i widma referencyjnego
\end{itemize}
